\documentclass[10pt,twocolumn]{article}
\usepackage[T1]{fontenc}
\usepackage[utf8]{inputenc}
\usepackage{lmodern}
\usepackage[margin=1.8cm,top=2.4cm,bottom=2cm,columnsep=0.6cm]{geometry}
\usepackage{fancyhdr}
\usepackage{titlesec}
\usepackage{booktabs}
\usepackage{amsmath,amssymb,amsfonts}
\usepackage{microtype}
\usepackage{xcolor}
\usepackage{mdframed}
\usepackage{parskip}
\usepackage{enumitem}
\usepackage{hyperref}

\definecolor{rulered}{RGB}{180,30,30}
\definecolor{boxgray}{RGB}{245,245,245}
\definecolor{keywordgray}{RGB}{100,100,100}

\pagestyle{fancy}
\fancyhf{}
\fancyhead[L]{\textsc{\small Claude Mag}}
\fancyhead[C]{\small\textit{Machine Learning Research Digest}}
\fancyhead[R]{\small 2026-11-18}
\fancyfoot[C]{\small\thepage}
\renewcommand{\headrulewidth}{0.8pt}

\titleformat{\section}{\normalsize\bfseries\color{rulered}}{}{0pt}{}%
  [\vspace{-4pt}\color{rulered}\rule{\columnwidth}{0.4pt}]
\titlespacing{\section}{0pt}{8pt}{4pt}

\hypersetup{colorlinks=true,urlcolor=rulered,linkcolor=black}

\begin{document}
\setlength{\parindent}{0pt}
\setlength{\parskip}{4pt}

{\small\color{keywordgray}\textbf{Keywords:}
  Vision-Language-Action \textbullet{} Diffusion VLM \textbullet{} Robot Manipulation \textbullet{} Unified Architecture}\par\vspace{4pt}

{\LARGE\bfseries\raggedright Concertina: Embodying Manipulation and Cognition in One Diffusion VLM}\par\vspace{4pt}
{\large\itshape\raggedright Inserting Action Tokens Natively into a Diffusion Backbone Achieves 98.15\% on LIBERO With No Adapters}\par\vspace{6pt}
{\small\textbf{By} Chengcheng Wang, Wei He, Wenhao Li, Hongguang Li, Kai Han, Jianyuan Guo, Chang Xu}\par\vspace{2pt}
{\small\color{keywordgray}arXiv:2611.07513 \textbullet{} cs.CV, cs.AI, cs.RO \textbullet{} Submitted 5 Nov 2026}
\par\vspace{2pt}\noindent{\color{rulered}\rule{\columnwidth}{0.6pt}}\vspace{6pt}

Wang~et~al.\ propose Concertina, a vision-language-action model that integrates robotic manipulation directly into a diffusion VLM backbone by inserting continuous action tokens into the same token sequence as language---with no action-boundary special tokens, embodiment-specific adapters, or soft prompts.
Concertina achieves 98.15\% average success on LIBERO, 85.9\% on the 10,030-task LIBERO-Plus benchmark, and a 4.142 average completed chain length on CALVIN ABC-D, outperforming prior VLA models across all three benchmarks.

\begin{mdframed}[backgroundcolor=boxgray,linecolor=rulered,linewidth=1pt,
    innerleftmargin=6pt,innerrightmargin=6pt,
    innertopmargin=6pt,innerbottommargin=6pt]
\textbf{\small AT A GLANCE}\par\vspace{4pt}
\begin{description}[leftmargin=0pt,labelsep=4pt,itemsep=2pt,parsep=0pt,
    font=\small\bfseries]
  \item[Problem:] VLA models separate perception and action into modules with an interface mismatch: the VLM was not trained for action generation, and the action module was not trained for scene understanding.
  \item[Why it matters:] Unified representations enable the same backbone to reason about and act on the world without representational seams.
  \item[Method:] Continuous action tokens are projected into the VLM's token space and generated by the same diffusion denoising process as language tokens; joint training.
  \item[Result:] LIBERO 98.15\%, LIBERO-Plus 85.9\%, CALVIN ABC-D chain 4.142.
  \item[Evidence:] LIBERO, LIBERO-Plus, CALVIN ABC-D; comparisons with OpenVLA, RoboVLMs, $\pi_0$.
\end{description}
\end{mdframed}

\section{The Big Picture}

Robot learning has converged on a paradigm: use a large pre-trained VLM for scene understanding and task grounding, then append a task-specific action module.
This paradigm works, but it carries an architectural cost: the VLM and the action module are optimized for different objectives, and their interface---typically a fixed-size feature vector or a cross-attention bottleneck---is a representational seam.

Concertina eliminates the seam.
By treating robotic actions as continuous tokens that the diffusion VLM generates alongside language tokens, the model learns a single representational substrate for seeing, reasoning, and acting.
There is no interface to design, no adapter to tune, and no embodiment-specific component to replace when switching robots.

The key enabling insight is that diffusion VLMs already generate \emph{continuous} outputs through iterative denoising---a process that maps naturally to the continuous, multi-step nature of robotic action trajectories.

\section{Why Existing Approaches Fall Short}

\textbf{Autoregressive VLAs} (OpenVLA) discretize actions into tokens for autoregressive generation.
Discretization loses precision for fine motor control, and the sequential generation order creates temporal inconsistencies in the action trajectory.

\textbf{Diffusion VLAs with separate action heads} ($\pi_0$, Diffusion Policy) append a diffusion head to a frozen or lightly fine-tuned VLM.
The VLM's representations were not trained to support diffusion-based action generation; the action head must learn to work around this, which limits generalization.

\textbf{Adapter-based VLAs} (RoboVLMs) add embodiment-specific LoRA adapters or action heads per robot type.
This scales poorly---each new embodiment requires new adapter training---and the adapters are not shared across tasks, so each task-embodiment combination is an independent fine-tuning problem.

\section{Core Mechanism}

Concertina has three components that jointly participate in the same diffusion forward pass:
\begin{enumerate}[leftmargin=1.5em,itemsep=2pt]
  \item \textbf{Visual encoder:} Encodes the camera image into patch features.
  \item \textbf{Language encoder:} Encodes the task description.
  \item \textbf{Joint denoising backbone:} A diffusion transformer processes visual features, language tokens, and action tokens in a unified sequence, generating all outputs by iterative denoising.
\end{enumerate}

No special tokens separate the action segment from the language segment.
The boundary is implicit in the position encodings of each token type.

\section{Technical Formulation}

\subsection*{Action Tokenization}

An action trajectory $\mathbf{a}_{1:H} \in \mathbb{R}^{H \times d_a}$ is mapped to continuous tokens via a lightweight linear projection:
\begin{equation}
\mathbf{e}^a_h = W_a \mathbf{a}_h + \mathbf{b}_a, \quad h = 1,\ldots,H
\label{eq:proj}
\end{equation}
where $W_a \in \mathbb{R}^{d \times d_a}$ and $\mathbf{b}_a \in \mathbb{R}^d$ are the only action-specific parameters in the entire model.
The projected tokens $\mathbf{e}^a_{1:H}$ are appended to the visual and language token sequence.

\subsection*{Joint Diffusion Denoising}

The full token sequence $\mathbf{z} = [\mathbf{z}^v;\, \mathbf{z}^l;\, \mathbf{e}^a]$ is treated as a single noisy observation.
At diffusion step $k$, Gaussian noise is added to the action tokens:
\begin{equation}
\mathbf{z}^{(k)} = [\mathbf{z}^v;\, \mathbf{z}^l;\, \mathbf{e}^a + \sigma_k \boldsymbol{\epsilon}], \quad \boldsymbol{\epsilon} \sim \mathcal{N}(\mathbf{0}, \mathbf{I})
\label{eq:noise}
\end{equation}
The denoising backbone $f_\theta$ predicts the clean action tokens:
\begin{equation}
\hat{\mathbf{e}}^a = f_\theta(\mathbf{z}^{(k)}, k)_{[\text{action positions}]}
\label{eq:denoise}
\end{equation}
The backbone attends to visual and language tokens to condition the action prediction, without any cross-attention bottleneck or adapter.

\subsection*{Training Objective}

The training loss is the standard diffusion denoising objective on action tokens:
\begin{equation}
\mathcal{L}(\theta, W_a) = \mathbb{E}_{k, \boldsymbol{\epsilon}}\!\left[\bigl\|\hat{\mathbf{e}}^a - \mathbf{e}^a\bigr\|^2\right]
\label{eq:loss}
\end{equation}
plus the standard VLM loss on language tokens.
Both are optimized jointly: the backbone's visual and language representations are shaped by both understanding and action objectives simultaneously.

\section{Learning or Inference Procedure}

\textbf{Training:}
Joint training mixes VLM data (image captioning, visual QA) with robot demonstration data at a 3:1 ratio.
The backbone is never frozen: unlike adapter-based approaches, Concertina updates all parameters to integrate manipulation and cognition into the same representational space.
Training runs for 200K steps on 8 GPUs with batch size 256.

\textbf{Inference:}
Given a camera observation and task description, the model runs $K=10$ denoising steps starting from Gaussian noise in the action token positions.
After denoising, the action tokens are decoded back to motor commands via $W_a^\top$ (transposed projection).
The same model runs all tasks without embodiment-specific switches.

\section{What the Guarantee Says}

No formal guarantee is provided.
The empirical evidence covers three well-established manipulation benchmarks with a total of over 10,000 task variations in LIBERO-Plus.
The near-zero failure rate on standard LIBERO (98.15\%) and competitive CALVIN chain length (4.142 vs.\ $\pi_0$'s 3.41) provide strong evidence of generalization.

\section{Experimental Findings}

\begin{center}
\small
\begin{tabular}{lrrr}
\toprule
\textbf{Method} & \textbf{LIB.} & \textbf{LIB.+} & \textbf{CALVIN} \\
\midrule
OpenVLA          & 73.1\% & 42.3\% & 2.69 \\
RoboVLMs         & 81.4\% & 56.2\% & 3.14 \\
$\pi_0$          & 90.3\% & 72.1\% & 3.41 \\
Concertina       & \textbf{98.15\%} & \textbf{85.9\%} & \textbf{4.142} \\
\bottomrule
\end{tabular}
\end{center}

The LIBERO-Plus result (85.9\% across 10,030 tasks) is particularly significant: it reflects generalization to tasks not seen during training, and the gap over $\pi_0$ (13.8 pp) is larger than on standard LIBERO, suggesting that the unified architecture generalizes better when task diversity increases.

\section{Ablations and Interpretation}

\begin{itemize}[leftmargin=1em,itemsep=2pt]
  \item \textbf{Joint vs.\ frozen backbone:} Freezing the backbone during action training reduces LIBERO-Plus success from 85.9\% to 71.3\%, confirming that full joint training is necessary for the backbone to learn action-supporting representations.
  \item \textbf{Action token position:} Placing action tokens before language tokens reduces performance; appending them after language tokens (current design) performs best.
  \item \textbf{Projection dimension $d_a$:} $d_a = 7$ (joint angles) with $d = 1024$ gives optimal performance; larger $d_a$ (e.g., 24 for a full-body humanoid) degrades slightly, suggesting the projection may need scaling for higher-dimensional embodiments.
  \item \textbf{Denoising steps $K$:} $K=10$ is sufficient; $K=5$ degrades CALVIN chain length by 0.3; $K=20$ shows no improvement.
  \item \textbf{Data mixture ratio:} The 3:1 VLM-to-robot data ratio is optimal; reducing VLM data below 50\% of the mixture causes the backbone to lose visual grounding capabilities, reducing task description following accuracy.
\end{itemize}

\section*{Reference}
C.\ Wang, W.\ He, W.\ Li, H.\ Li, K.\ Han, J.\ Guo, C.\ Xu.
\textbf{Concertina: Embodying Manipulation and Cognition in One Diffusion Vision-Language Model.}
arXiv:2611.07513. cs.CV, cs.AI, cs.RO. November 2026.
\url{https://arxiv.org/abs/2611.07513}

\end{document}
